\chapter{Problem Analysis and System Requirements}
\label{ch:requirements}

\section{Existing Problem or Process}
\placeholder{Describe the process that existed before the proposed solution. Include inputs, human or system actions, outputs, delays, and known failure points where evidence is available.}

\section{Problem Analysis}
\placeholder{Analyze the root causes and consequences of the problem. Support the analysis with observations, interviews, process evidence, or project records where permitted.}

\section{Proposed Solution}
\placeholder{Describe the proposed SMS-based fraud detection system at a high level. Explain the intended output and how it supports, rather than overstates, decision-making.}

\section{System Users and Actors}
\placeholder{Identify system users, external systems, administrators, analysts, or other actors. Remove actors that do not apply.}

\section{Functional Requirements}
\begin{longtable}{p{0.18\textwidth}p{0.28\textwidth}p{0.46\textwidth}}
  \caption{Functional requirements template}\label{tab:functional-requirements}\\
  \toprule
  \textbf{ID} & \textbf{Requirement} & \textbf{Acceptance or verification condition}\\
  \midrule
  \endfirsthead
  \toprule
  \textbf{ID} & \textbf{Requirement} & \textbf{Acceptance or verification condition}\\
  \midrule
  \endhead
  FR-01 & \placeholder{Describe input handling.} & \placeholder{State how it will be verified.}\\
  FR-02 & \placeholder{Describe preprocessing or normalization.} & \placeholder{State how it will be verified.}\\
  FR-03 & \placeholder{Describe classification or scoring.} & \placeholder{State how it will be verified.}\\
  FR-04 & \placeholder{Describe result presentation or storage.} & \placeholder{State how it will be verified.}\\
  FR-05 & \placeholder{Add another project-specific requirement.} & \placeholder{State how it will be verified.}\\
  \bottomrule
\end{longtable}

\section{Non-Functional Requirements}
\begin{longtable}{p{0.22\textwidth}p{0.70\textwidth}}
  \caption{Non-functional requirements template}\label{tab:nonfunctional-requirements}\\
  \toprule
  \textbf{Category} & \textbf{Requirement and verification method}\\
  \midrule
  \endfirsthead
  \toprule
  \textbf{Category} & \textbf{Requirement and verification method}\\
  \midrule
  \endhead
  Performance & \placeholder{Specify response-time, throughput, or resource requirements only if known.}\\
  Reliability & \placeholder{Specify expected behavior under errors or unavailable dependencies.}\\
  Usability & \placeholder{Specify requirements for users interpreting detection results.}\\
  Maintainability & \placeholder{Specify code, configuration, or model maintenance requirements.}\\
  Security & \placeholder{Specify access-control, logging, or secure-processing requirements.}\\
  Privacy & \placeholder{Specify handling, minimization, retention, and anonymization requirements.}\\
  \bottomrule
\end{longtable}

\section{System Constraints}
\placeholder{Document constraints involving data availability, compute resources, libraries, integration points, regulations, time, or organizational policies.}

\section{Assumptions}
\placeholder{List assumptions made during analysis and explain how each assumption could affect the results if it is false.}

\section{Data Requirements}
\placeholder{Describe required fields, formats, labels, provenance, volume, quality, and access conditions. Use placeholders for all quantities.}

\section{Security and Privacy Requirements}
\placeholder{Describe how sensitive message content, identifiers, credentials, logs, model artifacts, and access permissions must be protected. Mention applicable policies or regulations only when verified.}

\section{Fraud-Related Considerations}
\placeholder{Discuss class imbalance, changing fraud patterns, explainability, cost of missed fraud versus false alarms, and the distinction between model output and confirmed fraud.}

\section{Use Cases and Workflows}
\insertfigure{figures/use-case-diagram.pdf}{Placeholder for a use-case or workflow diagram}{fig:use-case}
\placeholder{Replace the placeholder figure with a diagram created from the actual requirements. Ensure actors, system boundaries, and data flows are accurate.}

\section{Chapter Summary}
\placeholder{Summarize the requirements and constraints that drive the design choices in \chapterref{ch:design}.}
